\section*{第 \ref{ch:g8:fractions} 章 --- 分数的四则运算}

\begin{solution}{exo:g8:fractions:1}
$\dfrac34 + \dfrac25 = \dfrac{15}{20} + \dfrac{8}{20} =
\dfrac{23}{20}$.

$\dfrac56 - \dfrac38 = \dfrac{20}{24} - \dfrac{9}{24} =
\dfrac{11}{24}$.

$\dfrac{7}{10} + \dfrac{8}{15} = \dfrac{21}{30} + \dfrac{16}{30} =
\dfrac{37}{30}$.
\end{solution}

\begin{solution}{exo:g8:fractions:2}
$-\dfrac27 + \dfrac{3}{14} = -\dfrac{4}{14} + \dfrac{3}{14} =
-\dfrac{1}{14}$.

$\dfrac16 - \dfrac54 = \dfrac{2}{12} - \dfrac{15}{12} =
-\dfrac{13}{12}$.

$-\dfrac35 - \dfrac12 = -\dfrac{6}{10} - \dfrac{5}{10} =
-\dfrac{11}{10}$.
\end{solution}

\begin{solution}{exo:g8:fractions:3}
$\dfrac58 \times \dfrac{4}{15} = \dfrac{5 \times 4}{8 \times 15}
= \dfrac{1}{2 \times 3} = \dfrac16$（约去 $5$ 与 $4$）。

$\dfrac{9}{14} \times \dfrac73 = \dfrac{9 \times 7}{14 \times 3}
= \dfrac{3}{2}$（约去 $7$ 与 $3$）。

$\left(-\dfrac65\right) \times \dfrac{10}{9} = -\dfrac{60}{45}
= -\dfrac43$（异号，结果为负）。
\end{solution}

\begin{solution}{exo:g8:fractions:4}
倒数：$\dfrac73$；\quad $\dfrac15$；\quad $4$；\quad
$-\dfrac92$（倒数保持原数的符号）。
\end{solution}

\begin{solution}{exo:g8:fractions:5}
$\dfrac49 \div \dfrac23 = \dfrac49 \times \dfrac32 = \dfrac{12}{18}
= \dfrac23$.

$\dfrac75 \div 14 = \dfrac75 \times \dfrac{1}{14} = \dfrac{7}{70}
= \dfrac{1}{10}$.

$6 \div \dfrac34 = 6 \times \dfrac43 = 8$.
\end{solution}

\begin{solution}{exo:g8:fractions:6}
$\dfrac{3/4}{9/8} = \dfrac34 \times \dfrac89 = \dfrac{24}{36} =
\dfrac23$.

$\dfrac{5/6}{10} = \dfrac56 \times \dfrac{1}{10} = \dfrac{5}{60} =
\dfrac{1}{12}$.

$\dfrac{2}{4/7} = 2 \times \dfrac74 = \dfrac72$.
\end{solution}

\begin{solution}{exo:g8:fractions:7}
先计算乘法：
$\dfrac13 + \dfrac12 \times \dfrac45 = \dfrac13 + \dfrac{4}{10}
= \dfrac{10}{30} + \dfrac{12}{30} = \dfrac{22}{30} = \dfrac{11}{15}$.

先计算括号：
$\left(\dfrac13 + \dfrac12\right) \times \dfrac45
= \dfrac56 \times \dfrac45 = \dfrac{20}{30} = \dfrac23$.
\end{solution}

\begin{solution}{exo:g8:fractions:8}
已装部分：$\dfrac25 + \dfrac13 = \dfrac{6}{15} + \dfrac{5}{15} =
\dfrac{11}{15}$（占整个水箱）。空的部分：
$1 - \dfrac{11}{15} = \dfrac{4}{15}$.
\end{solution}

\begin{solution}{exo:g8:fractions:9}
人数多于 $14$ 时的分数：
$\dfrac23 \times \dfrac34 = \dfrac{6}{12} = \dfrac12$，即全班人数的一半。
若这部分为 $12$ 人，则全班有 $24$ 人。
\end{solution}

\begin{solution}{exo:g8:fractions:10}
$\dfrac{15}{2} \div \dfrac34 = \dfrac{15}{2} \times \dfrac43
= \dfrac{60}{6} = 10$ 段 --- 结果为整数，绳子没有剩余。
\end{solution}

\begin{solution}{exo:g8:fractions:11}
$E$：先计算除法：
$\dfrac56 \div \dfrac{10}{9} = \dfrac56 \times \dfrac{9}{10}
= \dfrac{45}{60} = \dfrac34$；然后
$E = \dfrac23 - \dfrac34 = \dfrac{8}{12} - \dfrac{9}{12} =
-\dfrac{1}{12}$.

$F = \left(-\dfrac12\right)^2 \times \dfrac83 = \dfrac14 \times
\dfrac83 = \dfrac{8}{12} = \dfrac23$（平方结果为正）。
\end{solution}

\begin{solution}{exo:g8:fractions:12}
由内向外计算：$1 + 1 = 2$；然后
$1 + \dfrac12 = \dfrac32$；然后
\[
1 + \frac{1}{3/2} = 1 + \frac23 = \frac53 .
\]
\end{solution}

\begin{solution}{pb:g8:fractions:1}
\textbf{1.} $\frac12 + \frac13 = \frac36 + \frac26 = \frac56$ 且
$\frac12 + \frac14 = \frac24 + \frac14 = \frac34$。因此
$\frac56 = \frac12 + \frac13$ 与 $\frac34 = \frac12 + \frac14$
都是埃及分数（每式中的分数均为互不相同的单位分数）。

\textbf{2.} 取公分母 $28$：
\[
\frac14 + \frac{1}{28} = \frac{7}{28} + \frac{1}{28}
= \frac{8}{28} = \frac{2}{7} .
\]

\textbf{3.} $\frac13 + \frac16 = \frac26 + \frac16 = \frac36 =
\frac12$，且 $\frac14 + \frac{1}{12} = \frac{3}{12} +
\frac{1}{12} = \frac{4}{12} = \frac13$。一般地，
$\frac{1}{n+1}$ 与 $\frac{1}{n(n+1)}$ 的公分母为
$n(n+1)$：
\[
\frac{1}{n+1} + \frac{1}{n(n+1)}
= \frac{n}{n(n+1)} + \frac{1}{n(n+1)}
= \frac{n+1}{n(n+1)}
= \frac{1}{n} .
\]

\textbf{4.} 从 $1 = \frac12 + \frac12$ 出发，将
后一半按第 3 题拆分：
\[
1 = \frac12 + \frac13 + \frac16 ,
\]
三个互不相同的单位分数。

\textbf{5.} 将 $\frac13$ 在 $\frac56 = \frac12 +
\frac13$ 中拆分，得
\[
\frac56 = \frac12 + \frac14 + \frac{1}{12} .
\]
埃及分数的写法从来不是唯一的：拆分恒等式总能
把分母最大的单位分数替换为
两个新的、更小的单位分数——因此由任一写法
即可造出另一种写法。

\textbf{6.} 取分母为 $21$：$\frac27 = \frac{6}{21} <
\frac{7}{21} = \frac13$。取分母为 $28$：$\frac14 =
\frac{7}{28} < \frac{8}{28} = \frac27$。单位分数随分母增大而变小，
即 $1 > \frac12 > \frac13 > \frac14 >
\dots$；由于 $\frac13$（及其以上的所有数）都太大，而
$\frac14$ 合适，$\frac14$ 是小于
$\frac27$ 的最大单位分数。

\textbf{7.} $\frac27 - \frac14 = \frac{8}{28} - \frac{7}{28} =
\frac{1}{28}$：贪心法得到
$\frac27 = \frac14 + \frac{1}{28}$，恰为
第 2 题中书吏的记载。

\textbf{8.} 取分母为 $14$：$\frac12 = \frac{7}{14} <
\frac{10}{14} = \frac57$，而比它更大的唯一单位分数
$1$ 超过 $\frac57$。因此 $\frac12$ 是符合条件的最大者，
且
\[
\frac57 - \frac12 = \frac{10}{14} - \frac{7}{14} = \frac{3}{14} .
\]

\textbf{9.} 取分母为 $28$：$\frac{3}{14} = \frac{6}{28} <
\frac{7}{28} = \frac14$，故 $\frac14$ 太大。取
分母为 $70$：$\frac15 = \frac{14}{70} < \frac{15}{70} =
\frac{3}{14}$，故 $\frac15$ 合适：它是
小于 $\frac{3}{14}$ 的最大单位分数。于是
\[
\frac{3}{14} - \frac15 = \frac{15}{70} - \frac{14}{70}
= \frac{1}{70},
\qquad\text{因此}\qquad
\frac57 = \frac12 + \frac15 + \frac{1}{70} .
\]
验算：$\frac{35}{70} + \frac{14}{70} + \frac{1}{70} =
\frac{50}{70} = \frac57$。

\textbf{10.} 第一步：$\frac12 = \frac{5}{10} < \frac{8}{10} =
\frac45$，且 $1 > \frac45$，故减去 $\frac12$：
$\frac45 - \frac12 = \frac{8}{10} - \frac{5}{10} = \frac{3}{10}$。
第二步：$\frac13 = \frac{10}{30} > \frac{9}{30} =
\frac{3}{10}$ 太大，而 $\frac14 = \frac{5}{20} <
\frac{6}{20} = \frac{3}{10}$ 合适，故减去 $\frac14$：
$\frac{3}{10} - \frac14 = \frac{6}{20} - \frac{5}{20} =
\frac{1}{20}$。结论：
\[
\frac45 = \frac12 + \frac14 + \frac{1}{20},
\qquad
\frac{10}{20} + \frac{5}{20} + \frac{1}{20} = \frac{16}{20}
= \frac45 . \checkmark
\]

\textbf{11.} 各次余数的分子依次为：
对 $\frac27$ 依次为 $2$、$1$；对 $\frac57$ 依次为 $5$、$3$、
$1$；对 $\frac45$ 依次为 $4$、$3$、$1$。每组数
\emph{严格递减}。

\textbf{12.} $\frac ab$ 与 $\frac1n$ 的公分母
为 $b \times n$：
\[
\frac{a}{b} - \frac{1}{n}
= \frac{a \times n}{b \times n} - \frac{b}{b \times n}
= \frac{a \times n - b}{b \times n} .
\]

\textbf{13.} 以公分母 $b \times (n-1)$ 通分后，
不等式 $\frac{1}{n-1} > \frac ab$ 比较的是
分子 $b$ 与 $a \times (n - 1)$：
\[
b > a \times (n - 1) = a \times n - a ,
\]
两边同加 $a - b$，得 $a > a \times n - b$。因此
贪心法每进行一步后，余数的分子
（第 12 题）都是比前一个严格更小的整数。
整数不可能严格递减下去：至多 $a$ 步之后，
余数的分子为 $1$——即一个单位分数——或为
$0$，方法即告停止。

\textbf{14.} 取公分母 $18$：
\[
\frac12 + \frac13 + \frac19
= \frac{9}{18} + \frac{6}{18} + \frac{2}{18}
= \frac{17}{18} ,
\]
其和\emph{小于} $1$：遗嘱遗漏了 $\frac{1}{18}$ 的
驼群。这一差额正是全题的关键。在 $18$ 头驼中，
三份 $\frac12$、$\frac13$、$\frac19$ 恰为整数 $9$、
$6$ 与 $2$，三者之和为 $9 + 6 + 2 = 17$：恰好是整个驼群，
第十八头驼则送回邻居处。而且无人可以抱怨：遗嘱
承诺，在 $17$ 头驼中，各份分别为 $\frac{17}{2} = 8.5$，
再为 $\frac{17}{3} \approx 5.67$，再为 $\frac{17}{9} \approx 1.89$
头，而每个孩子所得都严格多于此数（分别为 $9$、$6$ 与 $2$）。
被遗漏的 $\frac{1}{18}$ 正是邻居那头驼
暂时填补的空缺。

\textbf{15.} 第 4 题给出了修正后的遗嘱：
\[
\frac12 + \frac13 + \frac16 = 1 .
\]
各份额为整数，须驼群的头数能被
$2$、$3$ 及 $6$ 整除——即能被 $6$ 整除。
满足条件的最小驼群为 $6$ 头：份额分别为 $3$、$2$ 与 $1$，且
无剩余。
\end{solution}
